\documentclass[10pt,a4paper]{amsart}
\usepackage[margin=19mm]{geometry}
\usepackage{amsmath,amssymb,mathtools}
\usepackage[scr=rsfs,cal=euler]{mathalfa}
\usepackage{xfrac}
\usepackage[unicode,hidelinks,bookmarks=false]{hyperref}
\newcommand{\ie}{i.e.\ }
\newcommand{\cf}{cf.\ }
\newcommand{\resp}{respectively\ }
\newcommand{\Subsection}{\subsection}
\providecommand{\nobreakdash}{\nobreak\hbox{-}}
\setlength{\emergencystretch}{2em}
\multlinegap0pt
\usepackage{bm}
\usepackage{amssymb}
\usepackage{mathtools}
\usepackage{algorithm}\usepackage{algpseudocode}
\usepackage{cancel}
\newtheorem{theorem}{Theorem}
\newcommand{\abs}[1]{\left|#1\right|}
\newcommand{\absd}[1]{|#1|}    %default
\newcommand*{\bbN}{\mathbb{N}}
\title{The Sample Complexity of Learning Lipschitz Operators with respect to Gaussian Measures}
\author{Ben Adcock, Michael Griebel and Gregor Maier*}
\date{Source: arXiv:2410.23440v4 (CC BY 4.0)}
\begin{document}
\maketitle

\noindent\emph{Source and licence.} The Sample Complexity of Learning Lipschitz Operators with respect to Gaussian Measures. Version arXiv:2410.23440v4 (CC BY 4.0), distributed by arXiv under the Creative Commons Attribution 4.0 International licence, http://creativecommons.org/licenses/by/4.0/, as recorded in the arXiv licence metadata for this exact version and shown on the abstract page https://arxiv.org/abs/2410.23440v4. Full licence notice: \texttt{licenses/arxiv-2410.23440-CC-BY-4.0.txt}.\par\smallskip

\noindent\textit{Editorial scope.} Counted unit: the article's theorem u\_pi(s+1) (source0.tex lines 1072--1094). The article's complete original proof of it is delivered with it (source0.tex lines 1096--1184, one \texttt{proof} environment of 89 lines). No mathematics is written, repaired or re-derived: the statement and the proof are the article's own lines, in the article's own notation, and no retained line is substituted or re-flowed.  Retained uncounted as the source-local context the proof needs: lines 952--955; lines 963--967; lines 969--973; lines 975--978.  Nothing was omitted from any retained span, so the package carries no omission record.  No retained line contains a citation, so the wrapper carries no bibliography and no bibliography entry.  No retained line contains a figure environment, an image-inclusion command or drawing code, so no image is delivered and the package declares no asset dependency.  Editorial formatting only, in addition: an AMS \texttt{amsart} preamble that restates the article's own declarations and macros used by the retained lines, namely the environments \texttt{theorem} on separate counters, together with the standard packages those lines need.  The retained environments print in this document as Theorem 1 in order of appearance, whereas the article prints the same items with the numbering of its own full document.  The complete native source of the article as distributed in the arXiv e-print for this exact version is bundled unchanged as \texttt{sources/167707/source0.tex}.
    \begin{equation}
    \label{eq: anisotropic TD index set}
               \prod_{i=1}^d \frac{1}{a_i i} \leq \absd{\Lambda_{d,\bm{a}}^{\textup{TD}}} \leq \prod_{i=1}^d \left(\frac{1}{a_i i} + 1\right).
    \end{equation}

\begin{equation}
\label{eq: s-term: definition of S_varepsilon}
    S(\varepsilon) := \left\{\bm{\gamma} \in \Gamma: u_{\bm{\gamma}}^{-2} \leq 1 + \frac{1}{\varepsilon^2} \right\}
    = \left\{\bm{\gamma} \in \Gamma: \sum_{i = 1}^{\infty} \frac{\gamma_i}{\lambda_{\bm{b}, i}} \leq \frac{1}{\varepsilon^2} \right\}
\end{equation}

\begin{equation}
\label{eq: s-term: effective dimension d}
    d(\varepsilon)
    := \min \left\{l \in \bbN_0: \lambda_{\bm{b}, l+1} < \varepsilon^2 \right\} \in \bbN_0.
\end{equation}

\begin{equation}
\label{eq: s-term: S_varepsilon isomorphic to anisotropic TD index set}
    S(\varepsilon) \overset{\simeq}{\longrightarrow} \Lambda_{d(\varepsilon), \bm{a}}^{\textup{TD}}, \quad \bm{\gamma} \mapsto (\gamma_1, \dots, \gamma_{d(\varepsilon)}),
\end{equation}

\begin{theorem}[Typical decays of \(u_{\bm{\pi(s + 1)}}\)]
\label{thm: s-term: upper bounds for u_pi(s+1)}
Let \(\alpha, \beta > 0\).
\begin{itemize}  
    \item [(a)] \emph{Algebraic spectral decay:} Let \(\lambda_{\bm{b}, i} = i^{-\alpha}\) for every \(i \in \bbN\). Then, for every \(\delta, \eta > 0\), there exists \(\bar{s} = \bar{s}(\alpha, \delta, \eta) \in \bbN\) such that for every \(s \geq \bar{s}\),
    \begin{equation}
    \label{eq: s-term: upper bound: algebraic}
        u_{\bm{\pi(s + 1)}} \leq \eta \log(s)^{-\frac{1}{2(1/\alpha + \delta)}}.
    \end{equation}

    \item [(b)] \emph{Exponential spectral decay:} Let \(\lambda_{\bm{b}, i} = e^{-\alpha i^{\beta}}\) for every \(i \in \bbN\). Then, for every \(\delta > 0\), there exists \(\bar{s} = \bar{s}(\alpha, \beta, \delta) \in \bbN\) such that for every \(s \geq \bar{s}\),
    \begin{equation}
    \label{eq: s-term: upper bound: exponential}
        u_{\bm{\pi(s + 1)}} \leq e^{- \frac{1}{2} \alpha^{\frac{1}{\beta + 1}} \left( \frac{\beta + 1}{\beta + \delta}  \right)^{\frac{1}{1 + 1/\beta}} \log(s)^{\frac{1}{1 + 1/\beta}}}.
    \end{equation}
    
    \item [(c)] \emph{Double exponential spectral decay:} Let \(\lambda_{\bm{b}, i} = e^{-e^{\alpha i}}\) for every \(i \in \bbN\). Then, for every \(\delta, \eta > 0\), there exists \(\bar{s} = \bar{s}(\alpha, \delta, \eta) \in \bbN\) such that for every \(s \geq \bar{s}\),
    \begin{equation}
    \label{eq: s-term: upper bound: double exponential}
        u_{\bm{\pi(s + 1)}} \leq e^{- \eta \log(s)^{\frac{1}{1 + \delta}}}.
    \end{equation}
\end{itemize}
\end{theorem}

\begin{proof}
    All rates can be derived by suitably choosing the parameter \(\varepsilon\) in the set \(S(\varepsilon)\), defined in~\eqref{eq: s-term: definition of S_varepsilon}.
    Suppose that \( \varepsilon \leq \sqrt{\lambda_{\bm{b}, 1}} \), so that \(\lambda_{\bm{b}, i} \geq \varepsilon^2\) for every \(i \in [d(\varepsilon)]\) by definition of \(d(\varepsilon)\) (see~\eqref{eq: s-term: effective dimension d}). By the relation~\eqref{eq: s-term: S_varepsilon isomorphic to anisotropic TD index set}, we may use the upper size bound in~\eqref{eq: anisotropic TD index set} to compute
    \begin{equation}
    \label{eq: upper size bound for S_varepsilon}
        \abs{S(\varepsilon)} \leq \prod_{i=1}^{d(\varepsilon)} \left(\frac{\lambda_{\bm{b},i}}{i \varepsilon^2} + 1\right) 
        \leq \prod_{i=1}^{d(\varepsilon)} \frac{\lambda_{\bm{b},i}}{\varepsilon^2} \left(\frac{1}{i} + 1\right) 
        \leq (2 \varepsilon^{-2})^{d(\varepsilon)} \prod_{i=1}^{d(\varepsilon)} \lambda_{\bm{b},i}.
    \end{equation}
    %In the second step we used the fact that \(\lambda_{\bm{b}, i} \geq \varepsilon^2\) for \(i \in [d(\varepsilon)]\) by definition of \(d(\varepsilon)\) (see~\eqref{eq: s-term: effective dimension d}). 
    For brevity, we write in the following \(d = d(\varepsilon)\).


    \emph{Case (a).}
    Using~\eqref{eq: upper size bound for S_varepsilon} as well as the Stirling type estimate \(d^d \leq e^d d!\), we obtain
    \begin{equation}
    \label{eq: upper size bound for S_varepsilon: algebraic}
        \abs{S(\varepsilon)} \leq (2\varepsilon^{-2})^d \prod_{i=1}^d i^{-\alpha}
        = (2\varepsilon^{-2})^d (d!)^{-\alpha}
        \leq (2\varepsilon^{-2})^d e^{\alpha d} d^{-\alpha d}.
    \end{equation}
    Let \(0 < \varepsilon \leq 2^{- \alpha / 2}\). By definition of \(d\), we have \(d \leq \varepsilon^{-2/\alpha} < d + 1 \leq 2d\). We take the logarithm on both sides in~\eqref{eq: upper size bound for S_varepsilon: algebraic} and compute
    \begin{align*}
        \log(\abs{S(\varepsilon)}) 
        &\leq d \log(2 \varepsilon^{-2}) + \alpha d - \alpha d \log(d) \\
        &\leq \varepsilon^{-2/\alpha} \log(2\varepsilon^{-2}) 
        + \alpha \varepsilon^{-2/\alpha} 
        - \frac{1}{2} \alpha \varepsilon^{-2/\alpha} \log\left( \frac{1}{2} \varepsilon^{-2/\alpha} \right) \\
        &= \frac{1}{2} \varepsilon^{-2/\alpha} \log(2 \varepsilon^{-2})
        + \left(\alpha + \frac{1}{2} \log(2^{\alpha + 1})\right) \varepsilon^{-2/\alpha}.
    \end{align*}
    Next, let \(\delta, \eta > 0\) be arbitrary. There exists \(0 < \bar{\varepsilon} = \bar{\varepsilon} (\alpha, \delta, \eta) \leq 2^{- \alpha / 2}\) such that for every \(0 < \varepsilon \leq \bar{\varepsilon}\), we have \(\log(2 \varepsilon^{-2}) \leq \eta^{2(1/\alpha + \delta)} \varepsilon^{-2\delta}\) as well as \(\alpha + \frac{1}{2} \log(2^{\alpha + 1}) \leq \frac{1}{2} \eta^{2(1/\alpha + \delta)} \varepsilon^{-2 \delta}\). It follows that
    \begin{equation}
    \label{eq: log-size bound for S_varepsilon: algebraic}
        \log(\abs{S(\varepsilon)}) 
        \leq \eta^{2(1/\alpha + \delta)} \varepsilon^{-2/\alpha - 2\delta},
        \quad \forall \varepsilon \leq \bar{\varepsilon}.
    \end{equation}
    We set the right-hand side equal to \(\log(s)\) and solve for \(\varepsilon^2\), which gives
    \begin{equation}
    \label{eq: s-term: equation for varepsilon and s: algebraic}
        \varepsilon^2 = \varepsilon(s)^2 = \eta^2 \log(s)^{-\frac{1}{1/\alpha + \delta}}.
    \end{equation}
    We can now choose \(\bar{s} = \bar{s}(\alpha, \delta, \eta) \in \bbN\) sufficiently large such that the right-hand side in~\eqref{eq: s-term: equation for varepsilon and s: algebraic} is smaller than \(\bar{\varepsilon}^2\) for every \(s \geq \bar{s}\). Then,~\eqref{eq: log-size bound for S_varepsilon: algebraic} holds with \(\varepsilon = \varepsilon(s)\) and we conclude \(\absd{S(\varepsilon(s))} \leq s\) for every \(s \geq \bar{s}\).
    Since \(S(\varepsilon) = \pi([\absd{S(\varepsilon)}])\), it follows that \(\bm{\pi(s + 1)} \not\in S(\varepsilon)\) and consequently, 
    \[
    u_{\bm{\pi(s + 1)}}^2 
    \leq (\varepsilon(s)^{-2} + 1)^{-1} 
    \leq \varepsilon(s)^2
    \leq \eta^2 \log(s)^{-\frac{1}{1/\alpha + \delta}},\quad \forall s \geq \bar{s}.
    \]
   
    \emph{Case (b).} By~\eqref{eq: upper size bound for S_varepsilon}, we find
    \begin{equation}
    \label{eq: upper size bound for S_varepsilon: super-exponential}
        \abs{S(\varepsilon)} \leq (2 \varepsilon^{-2})^d \prod_{i=1}^d e^{-\alpha i^{\beta}}
        = (2 \varepsilon^{-2})^d e^{-\alpha \sum_{i=1}^d i^{\beta}}.
    \end{equation}
    Let \(0 < \varepsilon \leq e^{-\alpha / 2}\). By definition of \(d\), we have \(d \leq \alpha^{-1/\beta} \log(\varepsilon^{-2})^{1/\beta} < d + 1\). We take the logarithm on both sides of~\eqref{eq: upper size bound for S_varepsilon: super-exponential} and compute
    \begin{align*}
        \log(\abs{S(\varepsilon)}) 
        &\leq d \log(2\varepsilon^{-2}) - \alpha \sum_{i=1}^d i^{\beta} 
        \leq d \log(2\varepsilon^{-2}) - \alpha \int_0^d t^{\beta} dt \\
        &\leq \alpha^{-1/\beta} \log(\varepsilon^{-2})^{1 + 1/\beta} 
        + \alpha^{-1/\beta} \log(2) \log(\varepsilon^{-2})^{1/\beta} \\
        &\hspace{30ex} - \frac{\alpha}{\beta + 1} ( \alpha^{-1/\beta} \log(\varepsilon^{-2})^{1/\beta} - 1 )^{\beta + 1}.
    \end{align*}
    Next, let \(\widetilde{\delta} > 0\) and \(0 < \eta < 1\) be arbitrary. There exists \(0 < \bar{\varepsilon} = \bar{\varepsilon}(\alpha, \beta, \widetilde{\delta}, \eta) \leq e^{-\alpha/2}\) such that \(\log(2) \leq \widetilde{\delta} \log(\varepsilon^{-2})\) and \(\alpha^{-1/\beta} \log(\varepsilon^{-2})^{1/\beta} - 1 \geq \eta^{\frac{1}{\beta + 1}} \alpha^{-1/\beta} \log(\varepsilon^{-2})^{1/\beta}\) for every \(0 < \varepsilon \leq \bar{\varepsilon}\). Hence,
    \[
    \log(\abs{S(\varepsilon)}) 
    \leq \alpha^{-1/\beta} \left( 1 + \widetilde{\delta} -  \frac{\eta}{\beta + 1} \right) \log(\varepsilon^{-2})^{1 + 1/\beta}, \quad \forall \varepsilon \leq \bar{\varepsilon}.
    \]
    We now set \(\delta = 1 + \widetilde{\delta} (\beta + 1) - \eta\) and conclude similarly as in case (a).
    
    \emph{Case (c).} By~\eqref{eq: upper size bound for S_varepsilon}, we find that
    \begin{equation}
    \label{eq: upper size bound for S_varepsilon: double exponential}
        \abs{S(\varepsilon)} \leq (2 \varepsilon^{-2})^d \prod_{i=1}^d e^{-e^{\alpha i}}
        = (2 \varepsilon^{-2})^d  e^{-\sum_{i=1}^d  e^{\alpha i}}.
    \end{equation}
    Let \(0 < \varepsilon \leq e^{-e^{\alpha} / 2}\). Then, by definition of \(d\), we have \(d \leq \log(\log( \varepsilon^{-2})) / \alpha < d + 1\). We take the logarithm on both sides in~\eqref{eq: upper size bound for S_varepsilon: double exponential} and compute
    \begin{equation}
    \label{eq: upper size bound for S_varepsilon: double exponential II}
        \log(\abs{S(\varepsilon)}) 
        \leq d \log(2 \varepsilon^{-2}) - \sum_{i=1}^d e^{\alpha i} 
        \leq \frac{1}{\alpha} \log(\log(\varepsilon^{-2})) \log(2\varepsilon^{-2}) - \frac{\log(\varepsilon^{-2}) - 1}{e^{\alpha} - 1}  + 1.
    \end{equation}
    Next, let \(\delta,\eta > 0\) be arbitrary. There exists \(0 < \bar{\varepsilon} = \bar{\varepsilon}(\alpha, \delta, \eta) \leq e^{-e^{\alpha} / 2}\) such that for every \(0 < \varepsilon \leq \bar{\varepsilon}\), the right-hand side in~\eqref{eq: upper size bound for S_varepsilon: double exponential II} is dominated by \((2 \eta)^{-(1 + \delta)} \log(\varepsilon^{-2})^{1+\delta}\). We can now conclude similarly as in case (a).
\end{proof}

\end{document}
